% from ega2-1-fr.tex
\oldpage[II]{125}

morphisme de $\mathcal{F}$,
$f^*(v\otimes1_{\mathcal{R}(Y)})$ est un endomorphisme de
$f^*(\mathcal{F}\otimes_{\mathcal{O}_Y}\mathcal{R}(Y))$, et on a un diagramme
commutatif
\[
\xymatrix{
f^*(\mathcal{F}\otimes_{\mathcal{O}_Y}\mathcal{R}(Y))
  \ar[r]^{f^*(v\otimes1)}\ar[d]_{\theta}
& f^*(\mathcal{F}\otimes_{\mathcal{O}_Y}\mathcal{R}(Y))\ar[d]^{\theta}\\
f^*(\mathcal{F})\otimes_{\mathcal{O}_X}\mathcal{R}(X)
  \ar[r]_{f^*(v)\otimes1}
& f^*(\mathcal{F})\otimes_{\mathcal{O}_X}\mathcal{R}(X).
}
\]
On en conclut aisément que $\det f^*(v)$ est l'image canonique, par
l'homomorphisme $f^*(\mathcal{R}(Y))\to\mathcal{R}(X)$, de la section
$\det v$ de $\mathcal{R}(Y)$~; on est en effet aussitôt ramené au cas où
$X=\operatorname{Spec}(A)$ et $Y=\operatorname{Spec}(B)$ sont affines, $A$ et
$B$ étant des anneaux intègres de corps de fractions respectifs $K$ et $L$,
l'homomorphisme $B\to A$ étant injectif et s'étendant donc en un monomorphisme
$L\to K$~; si $\mathcal{F}=\widetilde{M}$, où $M$ est un $B$-module de type
fini, le rang sur $L$ de $M\otimes_B L$ est \emph{égal} à celui de
$(M\otimes_B A)\otimes_A K$ sur $K$, et $\det((u\otimes1)\otimes1)$ est
l'image dans $K$ de $\det(u\otimes1)$ pour tout endomorphisme $u$ de $M$,
d'où la conclusion.
\end{env}

\begin{env}[6.4.10]
\label{II.6.4.10-fr}
Supposons enfin que $X$ soit un \emph{préschéma localement noethérien réduit},
de sorte que le faisceau $\mathcal{R}(X)$ des fonctions rationnelles sur $X$
est encore un $\mathcal{O}_X$-Module quasi-cohérent
(\hyperref[I.7.3.4-fr]{I, 7.3.4})~; soit d'autre part $\mathcal{E}$ un
$\mathcal{O}_X$-Module \emph{cohérent} tel que
$\mathcal{E}'=\mathcal{E}\otimes_{\mathcal{O}_X}\mathcal{R}(X)$ soit
\emph{localement libre} (de rang fini). En vertu de
(\hyperref[II.6.4.7-fr]{6.4.7}), si $\mathcal{E}'$ est de rang constant, on
peut, pour tout endomorphisme $u$ de $\mathcal{E}$, définir les $\sigma_j(u)$,
qui sont des sections de $\mathcal{R}(X)$ au-dessus de $X$. Lorsqu'on ne
suppose plus $\mathcal{E}'$ de rang constant, on peut encore définir
l'homomorphisme
$\det:\mathcal{H}om_{\mathcal{O}_X}(\mathcal{E},\mathcal{E})
\to\mathcal{R}(X)$.
\end{env}

\subsection{Norme d'un faisceau inversible.}
\label{subsection:II.6.5-fr}

\begin{env}[6.5.1]
\label{II.6.5.1-fr}
Soit $(X,\mathcal{A})$ un espace annelé et soit $\mathcal{B}$ une
$\mathcal{A}$-Algèbre (commutative). Le $\mathcal{A}$-Module $\mathcal{B}$
s'identifie canoniquement à un sous-$\mathcal{A}$-Module de
$\mathcal{H}om_{\mathcal{A}}(\mathcal{B},\mathcal{B})$, une section $f$ de
$\mathcal{B}$ au-dessus d'un ouvert $U$ de $X$ s'identifiant à la
multiplication par cette section. Si $(X,\mathcal{A})$ et $\mathcal{B}$
vérifient l'une des conditions énumérées dans
(\hyperref[II.6.4.8-fr]{6.4.8}), (\hyperref[II.6.4.9-fr]{6.4.9}) ou
(\hyperref[II.6.4.10-fr]{6.4.10}), on peut donc définir $\det(f)$ (et dans
certains cas les $\sigma_j(f)$) qui sont des sections de $\mathcal{A}$ ou de
$\mathcal{R}(X)$ au-dessus de $U$, que nous appellerons la \emph{norme de
$f$} (resp. les \emph{fonctions symétriques élémentaires de $f$}) et noterons
$N_{\mathcal{B}/\mathcal{A}}(f)$. Nous allons supposer vérifiée \emph{l'une}
des conditions suivantes~:
\begin{enumerate}
\item[(I)] \emph{$\mathcal{B}$ est un $\mathcal{A}$-Module localement libre
(de rang fini).}
\item[(II)] \emph{$(X,\mathcal{A})$ est un préschéma localement noethérien
réduit, $\mathcal{B}$ est un $\mathcal{A}$-Module cohérent tel que
$\mathcal{B}\otimes_{\mathcal{A}}\mathcal{R}(X)$ soit un
$\mathcal{R}(X)$-Module localement libre, et pour toute section
$f\in\Gamma(U,\mathcal{B})$ au-dessus d'un ouvert $U\subset X$,
$N_{\mathcal{B}/\mathcal{A}}(f)$ est une section de $\mathcal{A}$ au-dessus
de $U$.}
\end{enumerate}

